\selectlanguage{english}
\chapter{Background and Related Work}
\label{chap:background}

\section{Text-Attributed Graphs}

A \textit{text-attributed graph} (TAG) is a graph $G = (V, E)$ in which each node
$v \in V$ is associated with a text document $d_v$. TAGs are ubiquitous in practice:
academic citation networks attach paper abstracts to nodes, e-commerce product graphs
attach item descriptions, social networks attach user profiles, and encyclopaedic
knowledge graphs such as Wikipedia attach the full article text. The richness of the
node attributes distinguishes TAGs from ordinary relational graphs and opens the door
to combining graph-structural reasoning with natural language understanding.

Link prediction on TAGs therefore requires methods that can exploit both the topology
of $G$ and the semantics of $\{d_v\}$. Pure structural methods ignore the text
entirely and rely on the assumption that the neighbourhood structure alone provides
sufficient signal for deciding edge existence. This assumption is often violated in
sparse or evolving graphs, where a large fraction of node pairs share no common
neighbours even when a link genuinely exists or should exist. Conversely, pure text
methods ignore the graph entirely and may conflate conceptually related nodes (e.g.,
two articles on different battles in the same war) with linked ones. Hybrid methods
that combine both modalities have become the dominant paradigm, but they carry the
inference cost of both a graph encoder and a text encoder, making them challenging to
deploy at the scale of Wikipedia.

The DSAA 2023 Wikipedia benchmark \cite{phan2023link} is a canonical TAG link
prediction dataset: the graph has 838,000 nodes corresponding to Wikipedia articles,
and each node's text document is its Wikipedia description and optionally its infobox
content. This thesis treats the DSAA 2023 dataset as its primary evaluation arena and
considers the full spectrum from structural heuristics to Sentence-Transformer
encoding.

\section{Structural Methods}

\subsection{Heuristic Indices}

Structural heuristics quantify the topological proximity between two nodes using only
the adjacency structure of the graph. They are the oldest class of link prediction
methods and remain competitive baselines because they require no training data and are
computable in constant time per pair once neighbour lists are cached.

Let $\Gamma(u)$ denote the set of neighbours of node $u$ in $G$. The four heuristics
used in this thesis are defined as follows.

\textbf{Common Neighbours (CN)} counts the number of nodes that are adjacent to both
$u$ and $v$:
\begin{equation}
    \mathrm{CN}(u, v) = |\Gamma(u) \cap \Gamma(v)|.
    \label{eq:cn}
\end{equation}
The intuition is that two nodes are more likely to be linked if they already share many
mutual connections, analogous to the social principle that a friend of a friend is
likely a friend.

\textbf{Jaccard Coefficient} normalises the Common Neighbours count by the size of the
union of the two neighbourhoods, controlling for high-degree nodes:
\begin{equation}
    \mathrm{Jaccard}(u, v) = \frac{|\Gamma(u) \cap \Gamma(v)|}{|\Gamma(u) \cup \Gamma(v)|}.
    \label{eq:jaccard}
\end{equation}
A pair $(u, v)$ where $u$ is a hub with thousands of neighbours and $v$ has only two
neighbours will receive a lower Jaccard score than a pair of low-degree nodes sharing
the same number of common neighbours.

\textbf{Adamic-Adar Index} \cite{adamic2003friends} weights each common neighbour $w$
inversely by the logarithm of its degree, down-weighting hub nodes that are common to
almost every pair:
\begin{equation}
    \mathrm{AA}(u, v) = \sum_{w \in \Gamma(u) \cap \Gamma(v)} \frac{1}{\log |\Gamma(w)|}.
    \label{eq:aa}
\end{equation}
This index was originally proposed to measure the closeness of two people on the Web
by counting how many pages they both link to, discounting pages that are linked to by
nearly everyone.

\textbf{Preferential Attachment (PA)} \cite{barabasi1999emergence} is motivated by the
observation that in many real-world networks new edges preferentially attach to
high-degree nodes. The score is simply the product of the two node degrees:
\begin{equation}
    \mathrm{PA}(u, v) = |\Gamma(u)| \cdot |\Gamma(v)|.
    \label{eq:pa}
\end{equation}
Unlike the other three indices, PA does not require computing the intersection of
neighbour sets and is therefore computable in $O(1)$ once degrees are precomputed.

All four indices collapse to zero (for CN, Jaccard, and AA) or to a degree product
(for PA) when one or both nodes have no neighbours in the graph, which motivates the
cold-start handling described in Chapter~\ref{chap:methodology}.

\subsection{Graph Embeddings}

Graph embedding methods learn a low-dimensional continuous vector $\mathbf{z}_v \in
\mathbb{R}^d$ for each node $v$ such that pairs of nodes that are structurally
similar in $G$ have similar embeddings. Link prediction is then performed by combining
the embeddings of the two endpoint nodes, typically using the dot product, $L_1$
distance, or Hadamard product as the edge feature.

\textbf{Node2Vec} \cite{grover2016node2vec} is a random-walk-based embedding method
that generalises DeepWalk by introducing two hyperparameters $p$ and $q$ that control
the breadth-first search (BFS) versus depth-first search (DFS) bias of the random walk.
The \textit{return parameter} $p$ governs the probability of immediately returning to
the previously visited node: a low $p$ encourages the walk to stay local. The
\textit{in-out parameter} $q$ governs the probability of visiting nodes farther from
the walk's origin: a low $q$ encourages a DFS-like outward exploration, while a high
$q$ encourages a BFS-like sampling of the local neighbourhood. The transition
probability from node $v$ to node $x$, having just come from node $t$, is:
\begin{equation}
    \pi_{tx} = \alpha_{pq}(t, x) \cdot w_{vx},
\end{equation}
where $w_{vx}$ is the edge weight and $\alpha_{pq}(t, x)$ is $1/p$ if $x = t$, $1$
if $x$ is a common neighbour of $t$ and $v$, and $1/q$ otherwise. The generated
random walks are treated as sentences and fed to the Word2Vec skip-gram model to
produce node embeddings. Node2Vec thus provides a flexible framework for capturing
both local community structure and global network topology, depending on the choice of
$p$ and $q$.

While graph embeddings improve over raw heuristics by encoding multi-hop neighbourhood
information, they share the fundamental limitation of ignoring node text. They also
require recomputing embeddings when the graph changes, which is costly for Wikipedia,
whose link structure evolves continuously.

\section{Semantic Methods}

\subsection{Text Representations}

Before the advent of pre-trained language models, text-based link prediction relied on
classical natural language processing representations. \textbf{TF-IDF} (term
frequency-inverse document frequency) represents each document as a sparse vector
where the weight of term $t$ in document $d$ is proportional to how often it appears
in $d$ and inversely proportional to how many documents in the corpus contain it. A
pair $(u, v)$ can be represented by the cosine similarity of their TF-IDF vectors, by
their concatenation, or by the element-wise product. TF-IDF is computationally
efficient but treats text as a bag of words, ignoring word order and morphological
relationships.

\textbf{Part-of-speech (POS) tag frequency features} represent a document by
counting the frequency of each grammatical category (noun, verb, adjective, adverb,
etc.) in the text. The intuition is that different types of Wikipedia articles have
characteristic POS distributions: articles about events are verb-heavy, articles about
places are noun-heavy with many proper nouns, and articles about abstract concepts may
have more adjectives and prepositions. A 72-dimensional POS feature vector formed by
concatenating the POS frequency histograms of the two endpoint nodes provides a
lightweight but surprisingly informative signal for link prediction on this benchmark,
as demonstrated by Tran et al.\ \cite{tran2023text}.

\subsection{Pre-trained Language Models}

\textbf{BERT} \cite{devlin2019bert} (Bidirectional Encoder Representations from
Transformers) introduced the paradigm of pre-training a deep Transformer on large
text corpora using masked language modelling (MLM) and next-sentence prediction (NSP).
In MLM, a random 15\% of tokens in the input are replaced with a \texttt{[MASK]}
token, and the model is trained to predict the original token from the bidirectional
context. In NSP, the model is trained to predict whether two sentences are consecutive
in the original corpus. This pre-training instills rich syntactic and semantic
knowledge that can be transferred to downstream tasks by fine-tuning on labelled data.
BERT's bidirectional attention mechanism, which allows every token to attend to all
other tokens in both directions, contrasts with unidirectional models such as GPT and
gives BERT an advantage on tasks that require understanding of the full context.

\textbf{RoBERTa} \cite{liu2019roberta} (Robustly Optimised BERT Pre-training Approach)
removed the NSP objective, used dynamic masking (so that different tokens are masked on
each training pass), trained on more data with larger batch sizes, and found that these
changes consistently improved downstream task performance. RoBERTa became the
standard strong baseline for sentence-level NLP tasks.

\textbf{Sentence-BERT} \cite{reimers2019sentence} addressed the limitation that
vanilla BERT is slow for semantic similarity tasks because it requires a cross-encoder
that processes both sentences jointly. Sentence-BERT introduces a siamese network
architecture in which the same BERT model encodes each sentence independently into a
fixed-size sentence embedding, and similarity is computed as the cosine distance
between the two embeddings. This bi-encoder design reduces inference time from
$O(n^2)$ (cross-encoder over all pairs) to $O(n)$ (encode each node once) plus
$O(n^2)$ cosine similarity computations, which is much cheaper when the embeddings can
be precomputed. Sentence-BERT is fine-tuned on natural language inference datasets
using a softmax loss over entailment, neutral, and contradiction classes, or using a
contrastive loss that brings semantically similar sentence pairs closer together in
embedding space.

\section{Hybrid Architectures}

Hybrid methods combine graph topology with textual content, aiming to exploit both
sources of signal simultaneously. Several architectures have been proposed for
text-attributed graphs specifically.

\textbf{LinkBERT} \cite{yasunaga2022linkbert} augments BERT's pre-training by
replacing NSP with a \textit{document relation prediction} (DRP) task: instead of
predicting whether two randomly sampled sentences are consecutive, the model is trained
to predict whether two documents are hyperlinked. The pre-training corpus is
constructed by sampling document pairs from Wikipedia's hyperlink graph. This allows
the model to learn cross-document relationships that are invisible to standard BERT
training. LinkBERT was shown to improve performance on knowledge-intensive question
answering and multi-hop reasoning tasks.

\textbf{GraphFormers} \cite{yang2021graphformers} nest a Transformer text encoder
inside a GNN message-passing loop. At each GNN layer, each node's text encoding is
updated not only by its own text but also by the text encodings of its neighbours,
allowing structural and semantic information to be fused at multiple levels of
abstraction. This architecture achieves strong performance on node classification and
link prediction tasks on academic citation graphs but requires access to the full
graph at inference time and is expensive to train.

\textbf{Patton} \cite{jin2023patton} proposes two pre-training strategies for
text-attributed graphs: network-contextualized masked language modelling, which
conditions the masked token prediction on neighbour text, and masked node prediction,
which predicts whether a node belongs to the neighbourhood of a given anchor node from
its text alone. Patton provides both a structural and a semantic pre-training signal,
and was evaluated on paper classification and link prediction on the OGB citation
network.

\section{The DSAA 2023 Wikipedia Benchmark}

The DSAA 2023 Wikipedia link prediction benchmark \cite{phan2023link} was introduced
as a shared task at the 10th IEEE International Conference on Data Science and Advanced
Analytics. The dataset is constructed from a snapshot of the English Wikipedia
hyperlink graph and is designed to test link prediction methods on a large-scale,
text-rich graph.

The dataset consists of three files. The \textbf{training file} (\texttt{train.csv})
contains 948,232 node pairs with binary labels: 46\% positive (edge exists) and 54\%
negative (no edge). The class imbalance is mild and does not require aggressive
resampling, though it motivates the use of Macro F1-score as the evaluation metric
rather than accuracy. The \textbf{node file} (\texttt{nodes.tsv}) contains one row per
Wikipedia node with columns for the node identifier, the article title, a short
description, and an optional infobox field. Of the 838,000 nodes in the dataset, 668,000
appear in the structural (adjacency) component of the graph; the remaining nodes have
no structural neighbours and constitute the cold-start population. The \textbf{test
file} (\texttt{test.csv}) contains 238,000 node pairs without labels; the evaluation
metric is Macro F1-score submitted to the competition leaderboard.

The two top-performing submissions at DSAA 2023 are the primary baselines for this
thesis. Phan et al.\ \cite{phan2023link} formulate link prediction as natural language
inference (NLI): they concatenate the title and description of both nodes into a
hypothesis-premise pair and fine-tune an XLM-RoBERTa model to predict entailment
(positive link) or neutral/contradiction (no link). This approach achieves a Macro
F1-score of 1.0 on the leaderboard, suggesting that self-loops and
structurally obvious pairs are sufficient to achieve perfect score when the evaluation
set contains a large fraction of trivially classifiable pairs. Tran et al.\
\cite{tran2023text} demonstrate that a simpler POS-tag frequency feature combined with
a Random Forest classifier achieves a Macro F1-score of 0.99999, providing an
important data point that deep semantic understanding is not strictly necessary for
this benchmark.

\section{Gaps in Existing Work}

The existing literature on this benchmark and on link prediction in text-attributed
graphs more broadly leaves several important gaps that this thesis addresses.

\textbf{Cost transparency.} Prior top-performing systems report accuracy but not
inference cost. A method that calls a Sentence-Transformer on every test pair is orders
of magnitude more expensive than one that uses heuristics for the majority of pairs,
yet the two may achieve the same aggregate F1-score. CascadeLP is designed to make
this trade-off explicit by measuring the fraction of pairs resolved at each tier and
the throughput of the full cascade.

\textbf{Cold-start analysis.} Both Phan et al.\ and Tran et al.\ report aggregate
scores without disaggregating performance on cold-start pairs (nodes with no structural
neighbours). As Kim et al.\ \cite{kim2024llms} observe, cold-start is the dominant
regime in dynamic graphs, and a system that handles the cold-start case poorly will
degrade significantly as the graph evolves. CascadeLP's Tier~2 is specifically
designed for the cold-start case and is evaluated separately on this subpopulation.

\textbf{Dataset characterisation.} The near-perfect scores reported for this benchmark
invite the question of whether the benchmark is saturated or whether it contains a
large fraction of trivially classifiable pairs. This thesis provides a systematic
characterisation, identifying 10,221 self-loops (all positive and predictable without
any model), analysing the distribution of Common Neighbour counts, and quantifying
the cold-start fraction. This characterisation contextualises the reported scores and
is essential for understanding where future progress on this benchmark can come from.

\textbf{Hybrid efficiency.} Existing hybrid architectures such as GraphFormers and
Patton are expensive to train and deploy. CascadeLP takes a different approach: rather
than building a monolithic model that combines all signals at training time, it
sequences simple models and escalates to more expensive ones only when necessary. This
design is interpretable, modular, and easy to update when new feature families become
available.
